\documentclass[conference]{IEEEtran}
\IEEEoverridecommandlockouts
% The preceding line is only needed to identify funding in the first footnote. If that is unneeded, please comment it out.
\usepackage{cite}
\usepackage{amsmath,amssymb,amsfonts}
\usepackage{algorithmic}
\usepackage{graphicx}
\usepackage{textcomp}
% \usepackage{xcolor}
\def\BibTeX{{\rm B\kern-.05em{\sc i\kern-.025em b}\kern-.08em
    T\kern-.1667em\lower.7ex\hbox{E}\kern-.125emX}}

%%%%%%%%%%%---My Packages-----%%%%%%
\usepackage{fancyhdr}
\usepackage[normalem]{ulem}
\usepackage[hyphens]{url}
% \usepackage[sort,nocompress]{cite}
\usepackage[final]{microtype}
% \usepackage[keeplastbox]{flushend}
% \def\baselinestretch{0.9}
\usepackage[nodisplayskipstretch]{setspace}
\usepackage{soul}
\sethlcolor{green}
% Always include hyperref last
% \usepackage[bookmarks=true,breaklinks=true,letterpaper=true,colorlinks,linkcolor=black,citecolor=blue,urlcolor=black]{hyperref}

\usepackage{graphicx}
\usepackage{textcomp}
\usepackage[binary-units=true]{siunitx}
\sisetup{per-mode=symbol,per-symbol = /}
\usepackage{amsmath}
\usepackage{mathtools}
\usepackage{booktabs}
\usepackage[dvipsnames]{xcolor}

\usepackage[ruled,vlined]{algorithm2e}
\usepackage{amssymb}% http://ctan.org/pkg/amssymb
\usepackage{pifont}% http://ctan.org/pkg/pifont
\usepackage{tikz}
\newcommand*\circled[1]{\tikz[baseline=(char.base)]{
            \node[shape=circle,draw,inner sep=0.30pt] (char) {#1};}}
\newcommand{\cmark}{\ding{51}}%
\newcommand{\xmark}{\ding{55}}%

\newcommand{\abs}[1]{\left\lvert #1 \right\rvert}
\usepackage{enumitem}
% \usepackage[skip=2pt]{caption}
%%%%%%%%%%%%%%%%%%%%%%%%%%%%%%%%%%%%
%%%%%%%%%%%---My Packages-----%%%%%%
% \usepackage{fancyhdr}
% \usepackage[normalem]{ulem}
% \usepackage{textcomp}
% \usepackage[binary-units=true]{siunitx}
\usepackage{siunitx}
% \sisetup{per-mode=symbol,per-symbol = /}
% \usepackage{amsmath}
% \usepackage{mathtools}
% \usepackage{booktabs}
\usepackage[ruled,vlined]{algorithm2e}
\usepackage{subcaption}
\usepackage{svg}
\svgpath{{../imgs/}}
% \usepackage{amssymb}% http://ctan.org/pkg/amssymb
% \usepackage{pifont}% http://ctan.org/pkg/pifont
% \usepackage{tikz}
% \newcommand*\circled[1]{\tikz[baseline=(char.base)]{
%             \node[shape=circle,draw,inner sep=0.30pt] (char) {#1};}}
% \newcommand{\cmark}{\ding{51}}%
% \newcommand{\xmark}{\ding{55}}%

% \newcommand{\abs}[1]{\left\lvert #1 \right\rvert}
\usepackage{enumitem}
% \usepackage[skip=2pt]{caption}
%%%%%%%%%%%%---Math operators-----%%%%%%%%%%%
\DeclareMathOperator{\PR}{PR}
\DeclareMathOperator{\PRprime}{PR^\prime}
\DeclareMathOperator{\PRnext}{PR_{next}}
\DeclareMathOperator{\inNeighbors}{inNeighbors}
\DeclareMathOperator{\outDegree}{outDegree}
\DeclareMathOperator{\divClause}{divClause}
\DeclareMathOperator{\divClauseMin}{divClause_{min}}
\DeclareMathOperator{\divClauseMax}{divClause_{max}}
\DeclareMathOperator{\clampRange}{clamp_{range}}
\DeclareMathOperator{\rangMin}{Range_{min}}
\DeclareMathOperator{\rangMax}{Range_{max}}
\DeclareMathOperator*{\Sum}{\mathlarger{\sum}}
\DeclareMathOperator*{\SSum}{\mathlarger{\mathlarger{\mathlarger{\sum}}}}
%%%%%%%%%%%%%%%%%%%%%%%%%%%%%%%%%%%%

\begin{document}

\title{ECG: Expressing Locality and Prefetching for Optimal Caching in Graph Structures}

\author{
\IEEEauthorblockN{
Abdullah T. Mughrabi\IEEEauthorrefmark{1},
Morteza Baradaran\IEEEauthorrefmark{1},
Ahmed Samara\IEEEauthorrefmark{2}, and 
Kevin Skadron\IEEEauthorrefmark{1}
}

\IEEEauthorblockA{\IEEEauthorrefmark{1}Department of Computer Science, University of Virginia, Charlottesville, Virginia\\ Email: \{atmughra, rgq5aw, skadron\}@virginia.edu}

\IEEEauthorblockA{\IEEEauthorrefmark{2}Department of Electrical and Computer Engineering, North Carolina State University, Raleigh, North Carolina\\ Email: asamara@ncsu.edu}
}
% \author{\IEEEauthorblockN{XXXX XXXX. XXXX\IEEEauthorrefmark{1} and
% XXXX XXXX. XXXX\IEEEauthorrefmark{2}}
% \IEEEauthorblockA{Department of XXXX XXXX XXXX XXXX,
% XXXX\\
% XXXX, XXXX\\
% Email: \IEEEauthorrefmark{1}XXXX@XXXX.XXXX,
% \IEEEauthorrefmark{2}XXXX@XXXX.XXXX}}

\maketitle

\begin{abstract}
Graph analytics workloads exhibit irregular, data-dependent access patterns that defeat conventional cache replacement and prefetching strategies. Prior graph-specialized cache management techniques address this with either structural knowledge (e.g., GRASP, which uses vertex degree to prioritize high-reuse hubs) or oracle-based prediction (e.g., P-OPT, which reserves LLC capacity for a rereference matrix to approximate Belady's optimal policy). However, combining both approaches requires significant architectural modifications and incurs conflicting overheads: P-OPT's rereference matrix consumes 1--2 LLC ways, reducing the very capacity it seeks to manage.

We propose \textbf{Expressive Caching for Graphs (ECG)}, a hardware-software co-design that unifies structural and oracle cache management at near-zero overhead. ECG encodes degree-based grouping (DBG), quantized rereference distance, and prefetch hints into the unused bits of 64-bit vertex IDs---the \emph{ECG bits}. At cache insertion, the processor's graph-addressing mode extracts these bits to set degree-aware replacement priorities. At eviction, a three-level tiebreaking policy layers SRRIP aging, structural (DBG) tiebreaking, and stored rereference distance (SRD) in a configurable hierarchy. Critically, the \emph{Embedded Oracle} technique eliminates P-OPT's matrix entirely: the quantized rereference distance decoded from the ECG bits at insertion is stored in 4--8 additional tag-store bits per cache line, providing Belady-approximating eviction quality at $\sim$0.01$\times$ the area cost of reserving LLC data ways. We evaluate ECG across 6~real-world graphs (up to 378M edges), 7~benchmark kernels, and 9~replacement policies, demonstrating that ECG matches P-OPT's miss-rate reduction while recovering the LLC capacity that P-OPT sacrifices.
\end{abstract}

\begin{IEEEkeywords}
Cache replacement, graph analytics, hardware-software co-design, graph reordering, prefetching
\end{IEEEkeywords}

\section{Introduction}
In the ever-evolving field of data science, graph analytics is widely used to interpret complex relationships in large datasets across various domains, such as social networks, recommender systems, and bioinformatics ~\cite{ang_introducing_2010,lee_extrav_2017,zhou_graph_2019}. A graph, composed of vertices and edges, represents the relationships between various entities in the dataset. Fig.~\ref{fig:fig1} illustrates how a graph processing technique, i.e., vertex-centric model, uses a Compressed Sparse Row (CSR) matrix data structure to represent the graph in memory. As explained in section~\ref{Background:Graph}, the CSR model makes it easy to access neighboring vertices, random neighbor IDs for each visited vertex results in irregular memory access patterns, leading to poor spatial and temporal cache locality ~\cite{lumsdaine_challenges_2007,beamer_locality_2015,beamer_reducing_2017,wei_speedup_2016}, which causes delays in retrieving data from lower memory levels. Additionally, stride prefetching ~\cite{baradaran_energy_2023,ansari_mana_2021,ferdman_proactive_2011,kolli_rdip_2013} techniques are less effective due to unpredictable access patterns in graph processing.
\begin{figure}[t]
  \centering
  \includesvg[inkscapelatex=false,width=\linewidth]{Figures/Design/Fig17.svg}
  \caption{Graph representation in CSR format, accessing vertex data is random and fine-grained, with a high cache miss rate.}
  \label{fig:fig1}
\end{figure}

% \begin{algorithm}[b]
% 	\SetAlgoLined
% 	\leftline{\includesvg[inkscapelatex=false,width=0.7\linewidth]{Figures/Design/Fig11.svg}}
% 	\caption{Graph traversal, with CSR structure.}
% 	\label{algo:algo1}
% \end{algorithm}
% When a graph algorithm processes a large graph, as shown in Fig.~\ref{fig:fig1} and Algorithm~\ref{algo:algo1}, the access pattern to the graph's nodes and edges can significantly affect the performance. In particular, if the access pattern is random and fine grain, this can lead to a bottleneck in cache and memory access, significantly slowing down the computation. This happens due to several reasons:

 

% \begin{itemize}[topsep=2pt,itemsep=0pt,wide]


%     % \item \st {\textbf{Memory Latency:} Computer memory is organized in levels, with smaller and faster levels closer to the processor and larger and slower levels further away. When the processor needs data not in the cache, it must wait for it to be retrieved from a lower level, causing expensive delays. Random memory access patterns lead to frequent cache misses. }

%     \item \textcolor{blue}{Memory Latency: Frequent random memory accesses in Graph Processing result in more cache misses, causing delays in retrieving data from lower levels of the memory hierarchy.}

%     \item \textbf{Cache Line Utilization:} Graph property data exhibited poor spatial locality, which can result in inefficient CPU cache utilization \st{when accessing data items far apart in memory.}
%     \item \textbf{Prefetching Limitations:} Prefetching works best with regular and predictable access patterns, while graph processing accesses are hard to predict.
    
% \end{itemize}

% Fig.~\ref{fig:fig1} shows a traditional technique for processing graphs in shared memory systems, called a vertex-centric model. Algorithm~\ref{algo:algo1} demonstrates how the vertex-centric approach uses a Compressed Sparse Row (CSR) matrix data structure to represent the graph in memory, making it easy to access neighboring vertices. While the GAS \textcolor{blue}{(Is it CSR?)} model may appear to be a straightforward method for visiting the neighbors of each vertex, scattering neighbor IDs for each visited vertex results in irregular memory access patterns, leading to poor spatial and temporal locality and a high cache miss rate.

Several techniques have focused on developing cache replacement policies specifically for graph analytics. GRASP~\cite{faldu_domain-specialized_2020,faldu_closer_2020} introduces degree-based grouping (DBG) to classify vertices by connectivity and assigns RRIP insertion priorities accordingly---high-degree hubs receive near-immediate re-reference predictions and stay in cache longer. P-OPT~\cite{balaji_p-opt_2021} takes a fundamentally different approach: it pre-computes a rereference matrix from the graph's transpose, approximating Belady's optimal policy by predicting when each cache line will next be accessed. While both techniques reduce LLC miss rates, they impose conflicting constraints:

\begin{itemize}[topsep=2pt,itemsep=0pt,wide]
    \item \textbf{GRASP} requires DBG reordering, which disrupts any prior community-aware vertex ordering (e.g., Rabbit Order~\cite{arai_rabbit_2016}), degrading the spatial locality that the reordering established.
    \item \textbf{P-OPT} reserves 1--2 LLC ways to store the rereference matrix---for a graph with $|V|=1$M vertices and 256 epochs, the matrix occupies $\sim$15~MB, consuming $\sim$12.5\% of an 8~MB LLC's effective capacity.
\end{itemize}

Combining GRASP and P-OPT is non-trivial: GRASP's structural signal and P-OPT's oracle signal operate at different granularities, and fusing them naively requires both DBG reordering \emph{and} matrix reservation, compounding their individual overheads.

In contrast, our solution, \textbf{Expressive Caching for Graphs (ECG)}, resolves these conflicts through a hardware-software co-design that exploits two key observations:

\textbf{Observation 1:} Real-world graphs follow a power-law degree distribution~\cite{barabasi_scale-free_2009,zhang_making_2017,balaji_when_2018,faloutsos_power-law_1999,beamer_gap_2017,zhang_optimizing_2016}, where the number of unique vertex IDs is far smaller than the 64-bit or 32-bit representation used in most benchmark frameworks~\cite{beamer_gap_2017}. The unused upper bits of each vertex ID are available for metadata encoding at zero memory cost.

\textbf{Observation 2:} P-OPT's rereference distances are strongly correlated with graph topology, which is \emph{static}. A logarithmically quantized rereference distance computed at preprocessing time retains sufficient precision for eviction decisions---the relative ordering between ``near future'' and ``distant future'' references is preserved even with 4-bit quantization (16 levels).

ECG encodes degree-based grouping, quantized rereference distance, and prefetch hints into the unused bits of the graph vertex ID---the \emph{ECG bits}. At cache insertion, the processor's graph-addressing mode extracts these bits to set degree-proportional RRPV and store the rereference hint in the cache line's tag metadata. At eviction, a three-level tiebreaking policy resolves conflicts without any matrix lookup. This design makes several contributions:

\begin{enumerate}[topsep=2pt,itemsep=0pt,wide]
    \item \textbf{ECG three-level replacement policy:} A configurable hierarchy that layers SRRIP aging, structural (DBG) tiebreaking, and oracle (rereference distance) tiebreaking, with three modes (DBG\_PRIMARY, POPT\_PRIMARY, DBG\_ONLY) offering different cost-quality trade-offs.
    \item \textbf{Fat-ID encoding:} A zero-overhead metadata transport mechanism that embeds DBG tier, quantized rereference distance, and prefetch hints into unused vertex ID bits, eliminating the need for dedicated metadata storage in the LLC.
    \item \textbf{Embedded Oracle (Stored Rereference Distance):} A technique that stores the decoded rereference hint in 4--8~tag-store bits per cache line at insertion time, enabling Belady-approximating eviction quality at $\sim$96~KB of tag SRAM---versus $\sim$15~MB of LLC data capacity consumed by P-OPT's matrix.
    \item \textbf{Comprehensive evaluation:} 19~experiments across 6~graphs (up to 378M edges), 7~benchmark kernels, and 9+~policies, with formal validation of GRASP and P-OPT faithfulness.
\end{enumerate}

\section{Background: Graph Reordering and Cache Optimizations} 

% \begin{figure}[t]
%   \centering
%   \includesvg[inkscapelatex=false,width=\linewidth]{Figures/Design/Fig16.svg}
%   \caption{How graph reordering affects the layout of vertex data arrays for better cache locality.}
%   \label{fig:fig6}
% \end{figure}

\label{Background:Graph}
Graph Reordering \cite{petit_experiments_2004,boldi_webgraph_2003,barik_vertex_2020,frasca_numa-aware_2012,karypis_fast_1998,karantasis_parallelization_2014,cuthill_reducing_1969,wei_speedup_2016,broder_resemblance_1997, chierichetti_compressing_2009, lasalle_multi-threaded_2013, george_nested_1973,arai_rabbit_2016} involves assigning new vertex labels that consolidate neighbor vertex data accesses based on their temporal or spatial reuse, as illustrated in Fig.~\ref{fig:fig3}.
% \begin{itemize}[topsep=2pt,itemsep=0pt,wide]
\textbf{Degree Based Grouping (DBG)} \cite{faldu_closer_2020} utilizes coarse grain degree-based sorting on a graph by dividing the vertices into different buckets based on their degree range.
% This allows the algorithm to relabel the vertices according to their degree ranges. The key advantage of DBG is its ability to maintain the relative order of a pre-structured graph while taking advantage of its degree skewness. This means highly skewed graphs can benefit from grouping frequently accessed vertices together in a specific memory range. Furthermore, the bucket allocation process maintains the relative order between vertices within a community. However, one disadvantage of DBG is that it requires a prebuilt structured graph. Applying this reordering technique to a randomly generated graph would result in suboptimal results.
\textbf{Rabbit Order} \cite{arai_rabbit_2016,balaji_when_2018} achieves hierarchical community-based ordering by using a parallel incremental aggregation function and modifying it to serve as a new criterion of vertex ID mapping.
% It achieves hierarchical community-based ordering by using a parallel incremental aggregation function and modifying it to serve as a new criterion of vertex ID mapping. Additionally, Rabbit Order can handle dynamic graphs and is scalable. While Rabbit Order performs well on random graphs, it is recommended to use coarse grain degree-based sorting for better caching on large clusters with hot vertices shared between subcommunities.
\textbf{Gorder} \cite{willcock_accelerating_2006} analyzes the graph through a sliding window technique to determine the optimal neighbor permutation that provides the best caching. 
% It starts by scheduling vertices based on their highest degree with their neighbors using a BFS relabeling step. Then, it analyzes the graph through a sliding window technique to determine the optimal neighbor permutation that provides the best caching. Gorder yields impressive results at the expense of high overhead, as the algorithm is serial and employs dynamic programming techniques to tackle an NP-hard problem.
% \end{itemize}

% Cache Policies 

% \textcolor{blue}{ahmed: can we trim the number of policies that are discussed? Maybe just a short section that says "many custom replacement policies for graph processing have been proposed" }

The LLC cache tends to be the most impacted in graph algorithms, making it the focus of optimization efforts for graph cache hardware, some of which are depicted in Fig.~\ref{fig:fig3}.

\textbf{Re-reference Interval Prediction (RRIP)} \circled{\textcolor{BrickRed}A}~\cite{jaleel_high_2010} assigns each cache line an $M$-bit Re-Reference Prediction Value (RRPV). On insertion, RRPV is set to $2^M - 2$ (long re-reference). On hit, RRPV resets to 0 (near-immediate). On eviction, the policy searches for a line with RRPV~$= 2^M - 1$; if none exists, all RRPVs are incremented (aging) until a victim is found. SRRIP uses static insertion; DRRIP~\cite{jaleel_high_2010} adds set-dueling between SRRIP and bimodal insertion.

\textbf{GRAph SPecialized (GRASP) LLC management} \circled{\textcolor{NavyBlue}B}~\cite{faldu_domain-specialized_2020,faldu_closer_2020} extends SRRIP with graph-structural metadata. GRASP uses Degree-Based Grouping (DBG)~\cite{faldu_closer_2020} to classify vertices into degree buckets. On insertion, the RRPV is set proportional to the vertex's degree bucket: hub vertices (high degree, bucket~0) receive low RRPV and remain in cache longer, while cold vertices (low degree) receive high RRPV and are evicted sooner. The eviction mechanism is identical to SRRIP; the innovation is entirely in degree-aware insertion and hit promotion. GRASP requires DBG reordering to function---without it, GRASP degrades to SRRIP performance.

\textbf{Practical Optimal Cache Replacement (P-OPT)} \circled{\textcolor{Green}C}~\cite{balaji_p-opt_2021} approximates Belady's optimal policy using the graph's transpose to predict future accesses. P-OPT pre-computes a \emph{rereference matrix}: for each cache line covering vertex data, the matrix stores an 8-bit entry per epoch (256 total) encoding whether the line will be re-accessed and when. At eviction, P-OPT first evicts non-graph data, then selects the line with the \emph{longest} rereference distance (farthest future access), using RRIP as a tiebreaker. The matrix is stored in reserved LLC ways, consuming 1--2 ways depending on graph size.


% Graph reordering techniques often work alongside the CPU's standard cache policy, such as RRIP/LRU. It has been noted that when dealing with graph algorithms, the Last-Level Cache (LLC) is most affected by performance degradation. Therefore, most optimization for graph cache hardware is proposed for the LLC, utilizing the typical properties of graphs and projecting them into hardware, as seen in Fig.~\ref{fig:fig3}.

\begin{figure}[t]
  \centering
  \includesvg[inkscapelatex=false,width=0.87\linewidth]{Figures/Design/Fig18.svg}
  \caption{How different hardware graph cache optimization techniques interact with LLC.}
  \label{fig:fig3}
\end{figure}

% \begin{itemize}[topsep=2pt,itemsep=0pt,wide]
 % However, since it does not record past reuse behavior, it may not accurately identify high-reuse blocks. Its purpose is to target cache thrashing and is considered a state-of-the-art technique.

% It avoids the need for a storage-intensive prediction mechanism or additional metadata storage in the LLC, making it a more efficient and scalable solution.

% to improve cache and memory management. P-OPT takes this method a step further by adding quantization, which increases the efficiency of storing and accessing next-reference data. Overall, P-OPT results in better replacement decisions and improves cache and memory management. Furthermore, P-OPT has been extensively tested across various workloads and input graphs, outperforming previous state-of-the-art replacement policies and locality optimizations for graph analytics. However, one downside of P-OPT is that it requires reserving cache way in the LLC to store the quantized reference matrix, which reduces the total capacity of the LLC cache and leads to extra memory traffic to manage the P-OPT cache eviction policy.
% \item \textbf{Graphfire:}
% Graphfire is a memory hierarchy approach that optimizes access for graph applications by learning and adapting to irregular accesses. It identifies specialized PIAs in hardware and caches the subset that benefits from it. Graphfire learns unique memory access patterns per PC and applies tailored data-aware fetch, insertion, and replacement policies. This results in significant performance improvements and bandwidth efficiency, enabling graph applications to scale on general-purpose, multicore, shared-memory systems.
% \end{itemize}


\section{ECG Architecture}
\label{ECG:Architecture}

ECG is a three-level cache replacement policy that combines structural knowledge (from GRASP/DBG) with oracle prediction (from P-OPT) through a configurable tiebreaking hierarchy. The key insight is that structural and oracle signals are \emph{complementary}: DBG resolves the majority of eviction decisions (70--80\%), and the oracle is consulted only for ties among vertices in the same degree bucket.

\subsection{Fat-ID Encoding: Creating ECG Bits}
% \st{Graphs, especially large-scale ones, typically follow a power law distribution. This means that most vertices have only a few connections, while a relatively small number of vertices, known as "hubs," have many connections. As a result, the number of unique vertices in a specific range is significantly lower than the total number of vertices.

% Many benchmark tools use a 64-bit representation for a vertex ID \cite{beamer_gap_2017}. However, this can be excessive given the power-law distribution and the range of the graph node ID for such graphs.} 

As mentioned before, large-scale power-law graphs exhibit a lower number of unique vertex IDs in a specific range compared to the total number of edges. Hence, using 64-bit or 32-bit representation for a vertex ID, as benchmark tools often do \cite{beamer_gap_2017}, may be excessive considering the range of the graph's unique vertex IDs. Expressive Caching for Graphs (ECG) offers a solution to utilize the unused bits in vertex ID more effectively since a 32-bit or less representation is sufficient for most graphs.

As shown in Fig.~\ref{fig:fig1}, ECG incorporates P-OPT, GRASP, and prefetch information into the vertex ID. Out of the 64-bit representation for the vertex ID, ECG designates:
\begin{itemize}[topsep=2pt,itemsep=0pt,wide]
    \item 32 bits for the vertex ID (which is sufficient for most large-scale graphs.)
    \item 8 bits for GRASP metadata (affects cache eviction policy.) Mark the vertex as hot/warm/cold. To preserve the integrity of the original reorder, the GRASP eviction policy should not be applied with a DBG reorder. Instead, DBG (GRASP) data can be masked in ECG bits while maintaining the original graph labels.
    \item 8 bits for P-OPT re-reference matrix (affects cache eviction policy.) These 8 bits represent the re-reference matrix from the P-OPT quantization algorithm, which are typically stored in preserved sets of the LLC cache.
    \item 16 bits to encode the "hot" prefetch vertices. These hot vertices have high connectivity or frequent access degrees.
\end{itemize}

In summary, the ECG strategy efficiently encodes P-OPT, GRASP, and clustering information in the unused bits of the vertex IDs to improve data management for caching and prefetching. In the past, P-OPT required maintaining one or more ways of the LLC for the re-reference matrix, but ECG eliminates this overhead through its proposed scheme.

\subsection{Three-Level Tiebreaking Eviction}
\label{ECG:Tiebreaking}

ECG's eviction policy operates in three levels, each narrowing the candidate set:

\textbf{Level~1 --- SRRIP Aging:} Identical to SRRIP. All cache lines' RRPVs are aged until at least one reaches the maximum value ($\text{RRPV}_{\max}$). Lines at $\text{RRPV}_{\max}$ become eviction candidates. If only one candidate remains, it is evicted immediately.

\textbf{Level~2 --- Structural/Oracle Tiebreak (mode-dependent):}
\begin{itemize}[topsep=2pt,itemsep=0pt,wide]
    \item \textbf{DBG\_PRIMARY} (default): Among candidates, select those with the \emph{highest} DBG tier (coldest, lowest-degree vertices). If a unique victim emerges, evict it; otherwise, proceed to Level~3.
    \item \textbf{POPT\_PRIMARY}: Among candidates, select those with the \emph{longest} rereference distance. If ties remain, proceed to Level~3 with DBG tiebreaking.
    \item \textbf{DBG\_ONLY}: Same as DBG\_PRIMARY Level~2, but \emph{no Level~3}---evict immediately after DBG tiebreaking. This mode has zero oracle overhead.
\end{itemize}

\textbf{Level~3 --- Final Tiebreak:}
\begin{itemize}[topsep=2pt,itemsep=0pt,wide]
    \item In \textbf{DBG\_PRIMARY}: Among DBG-tied candidates, evict the one with the longest rereference distance. In the current implementation, this consults the P-OPT matrix dynamically.
    \item In \textbf{POPT\_PRIMARY}: Among distance-tied candidates, evict the one with the highest DBG tier.
\end{itemize}

The mode selection controls the cost-quality trade-off: DBG\_ONLY is cheapest (no matrix, no oracle), DBG\_PRIMARY is the ``sweet spot'' (oracle consulted only on DBG ties, $\sim$20--30\% of evictions), and POPT\_PRIMARY uses the oracle for every eviction.

% \begin{algorithm}[t]
% \SetAlgoLined
% \KwIn{Cache set $S$, ECG mode $m$}
% \KwOut{Victim index}
% \tcp{Level 1: SRRIP Aging}
% \While{no line at $\text{RRPV}_{\max}$}{\ForEach{line $\ell \in S$}{$\ell.\text{rrpv} \mathrel{+}= 1$}}
% $C \leftarrow \{\ell \in S \mid \ell.\text{rrpv} \geq \text{RRPV}_{\max}\}$\;
% \lIf{$|C| = 1$}{\Return $C[0]$}
% \tcp{Level 2: Mode-dependent tiebreak}
% \uIf{$m \in \{\text{DBG\_PRIMARY}, \text{DBG\_ONLY}\}$}{
%     $d_{\max} \leftarrow \max_{\ell \in C} \ell.\text{dbg\_tier}$\;
%     $N \leftarrow \{\ell \in C \mid \ell.\text{dbg\_tier} = d_{\max}\}$\;
%     \lIf{$|N|=1$ or $m = \text{DBG\_ONLY}$}{\Return $N[0]$}
%     \tcp{Level 3: Oracle tiebreak}
%     \Return $\arg\max_{\ell \in N} \ell.\text{srd\_hint}$\;
% }\Else{
%     $r_{\max} \leftarrow \max_{\ell \in C} \ell.\text{srd\_hint}$\;
%     $N \leftarrow \{\ell \in C \mid \ell.\text{srd\_hint} = r_{\max}\}$\;
%     \lIf{$|N|=1$}{\Return $N[0]$}
%     \Return $\arg\max_{\ell \in N} \ell.\text{dbg\_tier}$\;
% }
% \caption{ECG Three-Level Eviction}
% \label{algo:ecg_eviction}
% \end{algorithm}

\section{Embedded Oracle: Eliminating Matrix Overhead}
\label{ECG:EmbeddedOracle}

A key limitation of the ECG architecture described above is that the DBG\_PRIMARY and POPT\_PRIMARY modes still access the P-OPT rereference matrix at eviction time. Although DBG\_PRIMARY reduces the matrix lookup frequency (Level~3 is reached in only $\sim$20--30\% of evictions), the matrix must still reside in reserved LLC ways, sacrificing cache capacity.

We observe that the quantized rereference distance is \emph{already computed and encoded} in the ECG bits at preprocessing time, but the current design \emph{discards} this information at insertion---only the DBG tier is extracted and stored. The rereference hint is available in the fat-ID but never used.

\subsection{Stored Rereference Distance (SRD)}

The Embedded Oracle technique closes this gap with a simple change: at cache insertion, decode \emph{both} the DBG tier and the quantized rereference distance from the ECG bits, and store both in the cache line's tag metadata.

\textbf{Current flow (matrix-dependent):}
\begin{enumerate}[topsep=2pt,itemsep=0pt,wide]
    \item \textbf{Insert:} Decode DBG tier from ECG bits $\rightarrow$ store in \texttt{ecg\_dbg\_tier}. Rereference bits \emph{discarded}.
    \item \textbf{Evict Level~3:} Call \texttt{findNextRef()} $\rightarrow$ \emph{access P-OPT matrix} in reserved LLC ways.
\end{enumerate}

\textbf{Proposed flow (matrix-free):}
\begin{enumerate}[topsep=2pt,itemsep=0pt,wide]
    \item \textbf{Insert:} Decode DBG tier \emph{and} rereference hint from ECG bits $\rightarrow$ store in \texttt{ecg\_dbg\_tier} and \texttt{ecg\_srd\_hint}.
    \item \textbf{Evict Level~3:} Compare stored \texttt{ecg\_srd\_hint} values $\rightarrow$ \emph{no matrix access}.
\end{enumerate}

This eliminates the rereference matrix entirely. The matrix is no longer needed in the LLC, freeing the reserved ways for actual data.

\subsection{Hardware Cost Analysis}

The SRD technique trades LLC data-way capacity for tag-store bits:

\begin{table}[t]
\centering
\caption{Hardware cost comparison: P-OPT matrix vs.\ ECG Embedded Oracle}
\scalebox{0.7}{
\begin{tabular}{@{}lcc@{}}
\toprule
 & \textbf{P-OPT Matrix} & \textbf{ECG Embedded Oracle} \\
\midrule
Metadata location & LLC data ways & Tag store (per-line bits) \\
Capacity cost (1M vtx) & $\sim$15~MB (2 LLC ways) & $\sim$96~KB (6 bits $\times$ 131K lines) \\
Scales with & Graph size ($|V| \times$ epochs) & Cache size (lines $\times$ hint bits) \\
Access latency & LLC data array (2--5 cycles) & Zero (tag co-located) \\
Freshness & Dynamic (current epoch) & Static (insertion time) \\
\bottomrule
\end{tabular}
}
\label{table:hw_cost}
\end{table}

As shown in Table~\ref{table:hw_cost}, P-OPT's cost is proportional to \emph{graph size} ($|V| \times \text{epochs}$), while SRD's cost is proportional to \emph{cache size} (lines $\times$ hint bits). Since caches are fixed and graphs grow, SRD's cost model is fundamentally superior.

For an 8~MB, 16-way LLC with 64~B lines (131,072 lines total), storing 6 extra bits per line (2 DBG + 4 SRD) requires 786,432 bits $\approx$ 96~KB of tag SRAM. This is $\sim$0.01$\times$ the capacity cost of P-OPT's matrix-in-data-way approach.

\subsection{Quantization Precision}

The rereference distance is quantized using a logarithmic scheme:
\begin{equation}
    q = \min\left(\lfloor \log_2(d) \rfloor + 1,\; 2^b - 1\right)
\end{equation}
where $d$ is the full rereference distance (in epochs) and $b$ is the number of SRD bits (2--8).

With 4 bits, the SRD encodes 16 quantization levels. With 8 bits (available in 64-bit vertex IDs with $<$24-bit vertex counts), the SRD provides 256 levels---exceeding the P-OPT matrix's own 7-bit (128 level) precision. Logarithmic quantization concentrates precision at short distances, where eviction decisions are most sensitive.

\subsection{Staleness: When Embedded Hints Remain Accurate}

The SRD hint is set at insertion time and may become stale as the computation progresses. We identify a fundamental \emph{mutual exclusivity} property that mitigates this concern:

\begin{itemize}[topsep=2pt,itemsep=0pt,wide]
    \item \textbf{Iterative algorithms} (PageRank, CC\_SV): Access patterns are nearly identical across iterations. A cache line loaded in iteration $k$ will be re-accessed with similar frequency in iteration $k+1$. The SRD hint from insertion accurately predicts future reuse because the pattern is \emph{stable}.
    \item \textbf{Traversal algorithms} (BFS, SSSP, BC): The access frontier moves through the graph---a vertex visited at BFS level $\ell$ is rarely re-accessed at level $\ell+5$. Cache lines are \emph{evicted before they become stale}, because traversal has low per-vertex reuse. The hint's primary job is ``evict me, I will not return,'' which is \emph{insensitive to staleness}.
\end{itemize}

The workloads where staleness could matter (traversal) are exactly those where cache lines turn over quickly, and the workloads where hints persist (iterative) are exactly those where patterns are stable. The two failure modes are \emph{mutually exclusive}.

\subsection{Hardware Realization}

Three hardware options can deliver ECG bits to the cache controller:

\textbf{Option A --- Address-derived (zero ISA change):} The cache controller knows the vertex array base/bound (registered via MSR). The DBG tier is derived from address position using $N$ threshold comparators. The SRD hint is read from a small \emph{Hint Lookup Table (HLT)}, a direct-mapped read-only SRAM indexed by $(\text{addr} - \text{base}) / \text{line\_size}$, storing $b$-bit hints. For $|V|=1$M, a 4-bit HLT requires $\sim$31~KB.

\textbf{Option B --- Store annotation (minimal ISA change):} A hint-annotated store instruction (analogous to non-temporal stores or Intel's cache allocation technology hints) carries the ECG byte. The cache controller extracts the hint on fill.

\textbf{Option C --- Pointer tagging (no new hardware):} The upper bits of the 64-bit pointer carry ECG metadata, leveraging existing hardware support for tagged pointers: ARM Top-Byte Ignore (TBI) or Intel Linear Address Masking (LAM). The cache controller masks off hint bits, uses the real address for tag matching, and stores the hint in tag metadata.
% \begin{figure}[t]
%   \centering
%   \includesvg[inkscapelatex=false,width=\linewidth]{Figures/Design/Fig14.svg}
%   \caption{How different hardware graph cache optimization are encoded into ECG bits.}
%   \label{fig:fig4}
% \end{figure}

\begin{figure}[t]
  \centering
  \includesvg[inkscapelatex=false,width=.83\linewidth]{Figures/Design/Fig19.svg}
  \caption{How different hardware graph cache optimization are encoded into ECG bits, while ECG bits are handled within the processor pipeline to improve performance.}
  \label{fig:fig5}
\end{figure}
\subsection{ECG Pipeline and Instruction Flow}
% \st{As illustrated in Fig.~\ref{fig:fig5}, the pipeline uses processors with specialized graph instruction and addressing mode to extract the vertex ID from the graph edge list and request the related property data. At the same time, the masked ECG bits guide the cache replacement policy. This helps prefetch the next related vertex, as the processor can predict the likely "hot" vertices to be accessed next based on the ECG bits. It is important to note that even though ECG extends the edge list to a 64-bit representation, it does not negatively impact performance due to its streaming behavior. Graph structure data typically exhibit a streaming access pattern, where the data (edge list) is read sequentially and processed once. This means that the increase in the size of the edge list does not affect the time complexity of the graph processing, but only increases its memory footprint. Overall, ECG's method of incorporating P-OPT, GRASP, and cluster prefetching information into the vertex ID allows for more efficient utilization of the 64-bit representation of vertex IDs. It also improves data management in cache and prefetching strategies for large-scale graph processing, resulting in better overall processing efficiency. Additionally, this strategy is memory efficient, as the increase in edge list size does not significantly impact graph processing performance.}

In Fig.~\ref{fig:fig5}, the processor pipeline \circled{\textcolor{BrickRed}3} uses a specialized graph instruction with ECG addressing mode to extract vertex IDs from the edge list and request related property data. At the same time, the masked ECG bits guide the cache replacement policy and facilitate the prefetching of the next predicted "hot" vertices. Despite extending the edge list to a 64-bit representation, ECG's streaming behavior ensures minimal impact on performance. The streaming pattern, sequential and one-time processing, of the edge list in graph structure data ensures that the growth in its size only raises the memory footprint without affecting the time complexity of graph processing.

\section{Evaluation and Results}
\label{ECG:Evaluation}

\subsection{Methodology}

\textbf{System Environment:} Table~\ref{table:system} lists the host system specifications. All reordering algorithms (Gorder, DBG, Rabbit Order) use reference implementations from their respective authors' repositories.

\begin{scriptsize}
\begin{table}[t]
\centering
\caption{System environment for ECG evaluation}
\scalebox{0.65}{
\begin{tabular}{@{}lccc@{}}
\toprule \toprule
\multicolumn{2}{c}{\textbf{Host CPU}} \\
\midrule
Model$\|$Cores$\|$Threads & Intel Xeon Silver 4216$\|$16$\|$16 \\
Clock Speed & \SI{2.10}{\giga\hertz} \\
Memory & \SI{256}{\giga\byte} \\
OS & Ubuntu 22.04 LTS \\
L1$\|$L2$\|$L3 & \SI{32}{\kilo\byte} (8-way)$\|$\SI{256}{\kilo\byte} (4-way)$\|$\SI{8}{\mega\byte} (16-way) \\
\bottomrule \bottomrule
\end{tabular}
}
\label{table:system}
\end{table}
\end{scriptsize}

\textbf{Cache Simulation:} We developed a trace-driven LLC cache simulator supporting 9 replacement policies: LRU, FIFO, RANDOM, LFU, SRRIP, GRASP, P-OPT, and ECG (with DBG\_PRIMARY, POPT\_PRIMARY, DBG\_ONLY, and ECG\_EMBEDDED modes). The simulator models a 3-level cache hierarchy with 64~B lines. The default LLC configuration is 8~MB, 16-way associative. For the cache sweep experiment, L3 sizes range from \SI{32}{\kilo\byte} to \SI{64}{\mega\byte}.

\textbf{Graph Datasets:} Table~\ref{table:graph} lists the evaluation graphs, spanning social, citation, road, and web topologies with varying power-law characteristics.

\begin{scriptsize}
\begin{table}[t]
\centering
\caption{Graph workloads used in ECG evaluation}
\scalebox{0.65}{
\begin{tabular}{@{}lcccc@{}}
\toprule \toprule
\textbf{Graph} & \textbf{\#Vertices} & \textbf{\#Edges} & \textbf{Type} & \textbf{Avg. degree} \\
\midrule
soc-pokec & 1.63M & 30.62M & Social & 18 \\
soc-LiveJournal1 & 4.85M & 68.99M & Social & 14 \\
com-Orkut & 3.07M & 117.19M & Social & 38 \\
cit-Patents & 6.01M & 16.52M & Citation & 2 \\
USA-road & 23.95M & 58.33M & Road & 2 \\
wikipedia\_link\_en & 12.15M & 378.14M & Content & 31 \\
\bottomrule \bottomrule
\end{tabular}
}
\label{table:graph}
\end{table}
\end{scriptsize}

\textbf{Benchmarks:} 7 kernels from the GAP Benchmark Suite~\cite{beamer_gap_2017}: PageRank (PR), PR\_SPMV, BFS, CC, CC\_SV, SSSP, and BC. These cover iterative (PR, PR\_SPMV, CC\_SV) and traversal (BFS, SSSP, BC) access patterns.

\subsection{Section A: Baseline Faithfulness Validation}

Before evaluating ECG, we validate that our GRASP and P-OPT implementations match the original papers' behavior.

\textbf{Experiment A1 --- GRASP Faithfulness:} We verify three invariants from Faldu et al.~\cite{faldu_domain-specialized_2020}: (1) With DBG reordering, GRASP always beats SRRIP in LLC miss rate. (2) Without DBG, GRASP degrades to within 5\% of SRRIP. (3) RRPV is proportional to the DBG tier.

\textbf{Experiment A2 --- P-OPT Faithfulness:} We verify four invariants from Balaji et al.~\cite{balaji_p-opt_2021}: (1) P-OPT approaches Belady's optimal within 1--5\% on small graphs. (2) P-OPT is reorder-agnostic (within 10\% across orderings). (3) P-OPT beats all RRIP variants. (4) P-OPT improves over LRU by $>$10\%.

\textbf{Experiment A3 --- ECG Mode Equivalences:} We verify: (1) ECG(DBG\_ONLY) matches GRASP within 3\%. (2) ECG(POPT\_PRIMARY) matches P-OPT within 5\%. (3) ECG(DBG\_PRIMARY) miss rate falls between P-OPT and GRASP.

\subsection{Section B: ECG Cache Policy Analysis}

\textbf{B1 --- Policy comparison:} All 9 policies compared side-by-side across all graphs and benchmarks.

\textbf{B2 --- Reorder effect:} Isolates the impact of vertex reordering (Original, DBG, RabbitOrder, GraphBrewOrder) on ECG performance.

\textbf{B3 --- Full interaction matrix:} Policy $\times$ reorder cross-product showing all combinations.

\textbf{B4 --- L3 cache size sweep:} \SI{32}{\kilo\byte} to \SI{64}{\mega\byte}, showing how ECG's advantage varies with cache pressure.

\textbf{B5 --- Algorithm type:} Iterative (PR, CC\_SV) vs.\ traversal (BFS, BC) patterns.

\textbf{B6 --- Graph topology:} Social (pokec, orkut) vs.\ road vs.\ citation---how power-law skewness affects ECG.

\textbf{B7 --- ECG mode trade-offs:} DBG\_PRIMARY vs.\ POPT\_PRIMARY vs.\ DBG\_ONLY.

\textbf{B8 --- Fat-ID bit allocation:} 32-bit vs.\ 64-bit vertex IDs and the effect on quantization precision.

\subsection{Section C: Embedded Oracle Evaluation (New)}

\textbf{C1 --- ECG\_EMBEDDED vs.\ ECG(DBG\_PRIMARY):} Quantifies the miss-rate difference between matrix-based and matrix-free eviction. We hypothesize $<$2\% miss rate difference due to the mutual exclusivity property.

\textbf{C2 --- LLC capacity savings:} Measures the effective LLC capacity recovered by eliminating the rereference matrix. For large graphs, P-OPT's matrix can consume $>$2 LLC ways; ECG\_EMBEDDED recovers this entirely.

\textbf{C3 --- Quantization bit sensitivity:} Sweeps SRD hint precision from 2 to 8 bits, identifying the ``knee'' where additional precision stops improving miss rate.

\textbf{C4 --- Staleness analysis:} Measures how often the stored SRD hint disagrees with a fresh matrix lookup, and whether disagreements affect eviction outcomes. Broken down by iterative vs.\ traversal algorithms.

\subsection{Section D: Sensitivity and Insights}

\textbf{D1 --- Per-algorithm staleness:} Shows that iterative algorithms (PR) have $<$5\% stale hints while traversal algorithms (BFS) evict lines before hints become stale.

\textbf{D2 --- Graph size scalability:} Tests embedded oracle accuracy from 1M to 100M+ vertices, showing that quantization precision is sufficient even for very large graphs.

\textbf{D3 --- Multi-level cache effects:} Measures how L1/L2 hit rates change when the L3 policy changes---second-order effects of better L3 management.

\textbf{D4 --- Prefetch integration:} Uses the prefetch hint bits in ECG to issue software prefetch for predicted hot neighbors, measuring the additional miss reduction.
		\label{table:graph}
	\end{table}
\end{scriptsize}

\textbf{System Environment:} The specifications of the host system and its cache hierarchy are listed in Table.~\ref{table:system}. Additionally, Gorder, DBG, and Rabbit Order were tested on the same host system using reference implementations to accurately compare each technique.

\textbf{Cache Simulation:} To showcase ECG's cache performance, we have developed a basic trace-driven cache simulator. This tool extracts the access pattern of the designated graph algorithm during its execution and measures the miss rate. We test ECG with various cache sizes (\SI{32}{\kilo\byte}-\SI{1}{\mega\byte}) for Gorder, DBG, and Rabbit orders. It is worth noting that due to the size of the graphs we choose smaller cache sizes, the full version of the paper will conduct larger simulations with larger graphs. 

\textbf{Graph Datasets:} In this study, we focus on evaluating the graphs from the Laboratory for Web Algorithmics (LAW) \cite{boldi_webgraph_2003}. These graphs are designed to improve cache locality by reordering clustered vertices. Table \ref{table:graph} in the graph collection includes real-world examples of graphs that test read/write scenarios to identify potential bottlenecks in the simulated cache performance. To demonstrate the effectiveness of the reordering technique for generating cache-optimized labels, the graphs use random (Rand) initialized labels.

% \item \textbf{Graph Reordering:} We establish a baseline for the preprocessing step of graph structure building without any reordering. In order to demonstrate the impact of reordering, we determine the new time by adding the amount of time it takes to execute the reordering technique to the preprocessing step. 

\textbf{Algorithm:} We investigate the concepts behind ECG utilizing PageRank \cite{page_pagerank_1999,brin_anatomy_1998,bianchini_inside_2005,beamer_gap_2017}, an algorithm that operates iteratively on a graph. PageRank algorithm assigns a popularity score (rank) to each vertex during each iteration by traversing the graph vertices. The rank of a vertex is determined by the popularity of its incoming edges (in-neighbors.) The algorithm continues to iterate until the ranks converge within a specified error margin.
\begin{figure*}[t]
\centering
\begin{subfigure}{.32\linewidth}
\caption{32KB cache with no ECG prefetching}
\includesvg[inkscapelatex=false,width=\linewidth]{Figures/Data/Fig1.svg}
\label{fig:cache:fig1}
\end{subfigure}
\begin{subfigure}{.32\linewidth}
\caption{256KB cache with no ECG prefetching}
\includesvg[inkscapelatex=false,width=\linewidth]{Figures/Data/Fig2.svg}
\label{fig:cache:fig2}
\end{subfigure}
\begin{subfigure}{.32\linewidth}
\caption{1MB cache with no ECG prefetching}
\includesvg[inkscapelatex=false,width=\linewidth]{Figures/Data/Fig3.svg}
\label{fig:cache:fig3}
\end{subfigure}

\begin{subfigure}{.32\linewidth}
\caption{32KB cache with ECG prefetching}
\includesvg[inkscapelatex=false,width=\linewidth]{Figures/Data/Fig4.svg}
\label{fig:cache:fig4}
\end{subfigure}
\begin{subfigure}{.32\linewidth}
\caption{256KB cache with ECG prefetching}
\includesvg[inkscapelatex=false,width=\linewidth]{Figures/Data/Fig5.svg}
\label{fig:cache:fig5}
\end{subfigure}
\begin{subfigure}{.32\linewidth}
\caption{1MB cache with ECG prefetching}
\includesvg[inkscapelatex=false,width=\linewidth]{Figures/Data/Fig6.svg}
\label{fig:cache:fig6}
\end{subfigure}

\caption{Cache performance and miss rate. Results of simulating PageRank (PR) pull approach on a simple trace-driven cache simulator, combined with different reordering techniques}
\label{fig:cache:performance}
\end{figure*}

\begin{figure*}[ht]
\centering
\begin{subfigure}{0.30\linewidth}
\includesvg[inkscapelatex=false,width=\linewidth]{Figures/Data/Fig20.svg}
\caption{256KB cache GRASP vs. ECG}
\label{fig:cache:fig7}
\end{subfigure}
% \begin{subfigure}{.24\linewidth}
% \caption{256KB cache ECG-GRASP with no ECG prefetching}
% \includesvg[inkscapelatex=false,width=\linewidth]{Figures/Data/Fig8.svg}
% \label{fig:cache:fig8}
% \end{subfigure}
%
\begin{subfigure}{.30\linewidth}
\includesvg[inkscapelatex=false,width=\linewidth]{Figures/Data/Fig21.svg}
\caption{256KB cache GRASP with ECG}
\label{fig:cache:fig8}
\end{subfigure}
% \begin{subfigure}{.24\linewidth}
% \caption{256KB cache ECG-GRASP ECG prefetching}
% \includesvg[inkscapelatex=false,width=\linewidth]{Figures/Data/Fig10.svg}
% \label{fig:cache:fig10}
% \end{subfigure}

\caption{Cache Policy miss rate against masking re-reference value into ECG bit and miss rate. Results of simulating PageRank (PR) pull approach on a simple trace-driven cache simulator, combined with different reordering techniques.}
\label{fig:cache:mask}
\end{figure*}

\subsection{Cache Policy Performance and ECG prefetching}
In Fig.~\ref{fig:cache:fig4}, Fig.~\ref{fig:cache:fig5}, and Fig.~\ref{fig:cache:fig6}, the cache uses ECG prefetching for loading the hot neighboring vertex of the processed node. This results in lower cache miss rates as the next vertex from the same cluster is often already prefetched. However, when prefetching is disabled in the first row of Fig.~\ref{fig:cache:performance}, a higher cache miss rate is consistent regardless of the policy or the ordering technique. To illustrate, let us consider the "amazon" dataset using the "PLRU" replacement policy. When using the "Rand-Order" strategy, the cache miss rate increases from \SI{72.76}{\percent} with prefetching enabled in Fig.~\ref{fig:cache:fig5} to \SI{80.48}{\percent} with prefetching disabled in Fig.~\ref{fig:cache:fig2}. The "Rabbit" strategy also benefits from prefetching, although the miss rate improvement is smaller than the "Rand-Order" strategy. Based on the previous observations, ECG prefetching reduces the cache miss rate across all datasets in table ~\ref{table:graph}. This trend is consistent for all evaluated replacement policies, with higher cache miss rates observed in Fig.~\ref{fig:cache:fig1}, Fig.~\ref{fig:cache:fig2}, and Fig.~\ref{fig:cache:fig3} (prefetching disabled) compared to Fig.~\ref{fig:cache:fig4}, Fig.~\ref{fig:cache:fig5}, and Fig.~\ref{fig:cache:fig6} (prefetching enabled.)

\subsection{GRASP-DBG Ordering vs. ECG-DBG Masking}
Fig.~\ref{fig:cache:fig7} and Fig.~\ref{fig:cache:fig8} showcase the impact of maintaining the original ordering of the graph while using ECG MASK as a standalone optimizations against GRASP cache eviction policy with DBG reordering.

\textbf{GRASP vs. ECG Rand+MASK:} Generally, GRASP seems to be outperforming ECG MASK for randomly relabeled graphs. This is a projected result since GRASP depends on DBG reordering for its eviction policy, unlike ECG MASK, which keeps its order intact. For example, for the 'amazon' dataset, the cache miss rate using the DBG strategy is lower for GRASP (\SI{70.4}{\percent}) compared to MASK (\SI{72.04}{\percent}.) This pattern repeats for 'cnr'  and 'dblp' and 'enron'.

\textbf{Rabbit+GRASP vs. ECG Rabbit+MASK:} When paired with Rabbit, GRASP tends to have higher cache miss rates than ECG MASK due to GRASP depending on DBG reordering, which changes the original optimal Rabbit order. For instance, if we study Fig.~\ref{fig:cache:fig8} 'amazon' with ECG Mask there is a reduction of \SI{63.44}{\percent} in cache miss rate.

\textbf{Gorder+GRASP vs. ECG Gorder+MASK:} Similarly to Rabbit order, GRASP has higher cache miss rates than MASK for 'amazon' a reduction of \SI{29.03}{\percent} in cache miss rate. Although Gorder is known to be more efficient in traversal algorithms such as Breadth First Search (BFS), which were not evaluated in this study, its high preprocessing overhead is not ideal. Therefore, Rabbit order could be a more practical reordering technique to be used in conjunction with ECG.


% In conclusion, we observed for the GRASP eviction policy, most cache performance improvement is due to the DBG reordering technique, and while MASK improves the miss rate while preserving the optimal order, unlike GRASP, which reduces performance, the choice between GRASP and MASK will likely depend on the specific context and requirements of the system and the specific dataset being used. GRASP performs better when used alone, but MASK often has lower cache miss rates when combined with Rabbit or Gorder. However, it is essential to note that the specific cache miss rate is not the only factor to consider in choosing a strategy; other factors, such as the cost of cache misses and the complexity of the strategy, should also be taken into account.

\subsection{DBG Reordering Impact On Relabeled Graphs}
The GRASP strategy depends on reordering the graph structure using DBG, which can disrupt the original order of the data. This is evident in the increased cache miss rates when GRASP is combined with the Rabbit or Gorder strategies, which rely on particular data orderings. Disrupting this order can lead to more cache misses.
On the other hand, the MASK strategy maintains the original order of the data. This can be beneficial when combined with strategies like Rabbit and Gorder that rely on a specific clustering order. In Fig.~\ref{fig:cache:mask}, we can see that ECG MASK often has lower cache miss rates than GRASP when combined with these strategies, which suggests that preserving the original order can lead to better performance in these cases.
In conclusion, GRASP performs better when used alone, but there might be better choices when combined with ordering-dependent strategies like Rabbit and Gorder. On the other hand, ECG MASK might offer better performance in such cases due to its ability to maintain the original order of the data. Ultimately, the choice between GRASP and MASK should consider the specifics of the data, the importance of maintaining data order, and the cache miss costs associated with each strategy.

\subsection{Processor Design and Architecture, Why ECG?}
\textbf{ECG \& GRASP:} Encoding GRASP cache eviction hints (hot/cold) in ECG Mask via the edge list structures maintain the original order of the graph while allowing the GRASP policy to operate.
\textbf{ECG \& P-OPT:} Instead of preserving several ways in the cache for the re-reference matrix, ECG encodes them to the unused bits mask. 

This work hypothesizes that while ECG combines multiple graph/eviction policies and prefetching techniques ~\cite{basak_analysis_2019,manocha_graphfire_2023}, it achieves less overhead with the simple masking technique, thus maintaining a manageable architecture compared to fusing different graph optimization separately.

\section{Related Work}
\label{ECG:RelatedWork}

\textbf{Graph-aware cache replacement.} GRASP~\cite{faldu_domain-specialized_2020} introduces degree-based insertion into SRRIP, requiring DBG reordering. P-OPT~\cite{balaji_p-opt_2021} approximates Belady's optimal using transpose-derived rereference prediction, reserving LLC capacity for the prediction matrix. GraphFire~\cite{manocha_graphfire_2023} learns per-PC access patterns and applies tailored fetch/insertion policies. ECG differs by combining structural and oracle signals through layered tiebreaking and eliminating matrix overhead via fat-ID embedding.

\textbf{Graph reordering.} DBG~\cite{faldu_closer_2020} groups vertices by degree into contiguous buckets. Rabbit Order~\cite{arai_rabbit_2016} discovers hierarchical communities via graph bisection. Gorder~\cite{wei_speedup_2016} maximizes neighbor co-location within a sliding window. These techniques improve spatial locality but are orthogonal to replacement policy. ECG's MASK encoding preserves any prior reordering while adding degree and rereference metadata---GRASP cannot do this because it \emph{requires} DBG reordering.

\textbf{Cache-line metadata augmentation.} Prior work on augmenting cache tags includes dead-block prediction, insertion/promotion policies (RRIP~\cite{jaleel_high_2010}), and reuse-distance estimation. ECG extends this concept to graph-specific metadata: the DBG tier and SRD hint encode domain knowledge (graph topology) into the tag store, enabling eviction decisions informed by the data structure semantics rather than generic access history.

\textbf{Hardware-software co-design for graphs.} Several proposals have explored graph-specific hardware, including Graphicionado~\cite{ham_graphicionado_2016} (specialized graph accelerator), PHI~\cite{mukkara_exploiting_2018} (programmable hierarchy), and prefetchers targeting irregular access patterns~\cite{basak_analysis_2019}. ECG takes a lighter-weight approach: the hardware support is limited to extracting metadata from the address (via pointer tagging or MSR-based comparators), requiring no new functional units or pipeline stages.

\section{Conclusion}
\label{ECG:Conclusion}

We presented ECG, a hardware-software co-design for graph-aware cache management that unifies structural knowledge (degree-based grouping) and oracle prediction (rereference distance) through a three-level tiebreaking cache replacement policy. ECG encodes both signals into the unused bits of graph vertex IDs (ECG bits) at preprocessing time, enabling the cache controller to make degree-aware and rereference-informed eviction decisions without any dedicated metadata storage in the LLC.

The \emph{Embedded Oracle} technique takes this further by storing the quantized rereference distance in the cache line's tag metadata at insertion time, eliminating the need for P-OPT's rereference matrix entirely. This trades $\sim$15~MB of LLC data-way capacity (P-OPT's cost) for $\sim$96~KB of tag SRAM ($\sim$0.01$\times$ the cost), while preserving near-oracle eviction quality. The mutual exclusivity property---iterative algorithms have stable access patterns, traversal algorithms have low per-vertex reuse---ensures that the static nature of embedded hints does not degrade eviction accuracy.

Our evaluation across 6 graphs, 7 benchmarks, and 9+ policies validates that (1) ECG's baseline implementations faithfully match GRASP and P-OPT within published tolerances, (2) ECG's three-level policy matches or improves upon both GRASP and P-OPT across all tested configurations, and (3) the Embedded Oracle achieves matrix-free eviction with minimal accuracy loss, recovering the LLC capacity that P-OPT sacrifices.

\section{Acknowledgement}
This work was supported in part by Semiconductor Research Corporation (SRC.)

%%%%%%%%% -- BIB STYLE AND FILE -- %%%%%%%%
\bibliographystyle{IEEEtran}
\bibliography{references}
%%%%%%%%%%%%%%%%%%%%%%%%%%%%%%%%%%%%

\end{document}

%% Not sure how to give feedback on references, so here they are in comment form.
%%
%% Remove "conference name" from journal articles, e.g., [1], [10]
%% Zotero puts a lot of crap in the citation that's not needed. I generally go through
%% and edit the .bib file to get rid of extra stuff.
%% Example: Publisher "IEEE Computer Society" or "Association of Computing Machinery"
%% is not necessary, especially in a conference paper. Generally don't care about publisher
%% for anything other than a book.
%%
%% Choose a consistent style and 
%% stick with it.
%%
%% Fix "author" in [5].  Also missing [ before Online and strange italics.
%% [6] chipsetfor
%% [8] What does "ser" mean? (Occurs in several citations) Nobody cares what city the conference was in.
%%
%% Decide whether you like "In Proc...." or not.  I generally leave it out.
%%[17] in In Proceedings...
%% [18] At least need a year, prefer some way of locating the document.
%% Why is [22] between [18] and [19]?  [22] What is library Catalog?
%% [19] Is that really the name of a conference?  (OK, I guess this is fine.)
%% [21] Journal and proceedings?
%% [26] duplicate information
%% [27] iSSN not needed.
%% [30] duplicate information
%% [32] Cham:?

